\section{Demonstrators Functionality and Architecture} \label{Demonstrators Functionality and Architecture}

\begin{figure}[H]
	\centering
	\includegraphics[width=1\linewidth]{../../../Images/Syde/HTA/DemonstratorFlowChart}
	\caption[Flowchart of Demonstrator Systems and their Interactions]{Flowchart of Demonstrator Systems and their Interactions}
	\label{fig:DemonstratorFlowChart}
\end{figure}
The demonstrator consists of three interconnected subsystems. At its core is the \ac{HTA} system, implemented within a trailer and designed in accordance with the previously described architecture. Complementing this setup is a dashboard, realized as a mobile application, which connects to the system via WiFi to provide monitoring and interaction capabilities. The trailer itself is pulled by a mobile system that applies the external force required for the \ac{HTA} module to operate. This mobile system is controlled independently through a dedicated remote control, ensuring flexible maneuverability. The flowchart presented in the following figure \ref{fig:DemonstratorFlowChart} illustrates the overall structure of the demonstrator. By analyzing each subsystem in detail, the interaction and data exchange between the hardware components can be clearly understood.

The External Force system is controlled by a Remote Control Unit \ref{FlowChartController}, which is equipped with two KY-023 Joystick modules \cite{JoyIT:2017}, each managing the speed and direction of one motor of the system. A SSD1315 display \cite{Seed:2025} provides overview of the real-time velocity of the demonstrator. Additionally, a power switch and a push button are included to initiate an ESP-NOW connection \cite[170]{Karuppusamy:2022} with the External Force system for wireless data exchange.This remote is as well built around a Hypatia board, serving as both the communication hub and the control unit for the connected hardware components.

\begin{figure}[H]
	\centering
	\includegraphics[width=1\linewidth]{../../../Images/Syde/HTA/FlowChartController}
	\caption[Flowchart of Remote Control System]{Flowchart of Remote Control System}
	\label{fig:FlowChartController}
\end{figure}

The External Force system applies tensile and compressive loads to the \ac{HTA} system by towing the wagon. In this system (see figure \ref{fig:FlowChartExtForce}) a second Hypatia board receives throttle input and controls two independently driven motors via an MOTODRIVER2, featuring a L298N H-bridge driver \cite{SBC-MotoDriver2:2025} and regulating both power and direction. Motor velocity feedback is provided by two LM393 \cite{Rajguru:2017} encoders. Similar to the main system, a switch and LED enable power control and status monitoring. The system is powered by four 1.5V batteries and a 5V power bank.

\begin{figure}[H]
	\centering
	\includegraphics[width=1\linewidth]{../../../Images/Syde/HTA/FlowChartExtForce}
	\caption[Flowchart of External Force System]{Flowchart of External Force System}
	\label{fig:FlowChartExtForce}
\end{figure}

The main system (see figure \ref{FlowChartHTA}) is the wagon including the \ac{HTA} system and a WiFi dashboard that provides graphical visualization of measurement data and enables real-time adjustment of control parameters. The \ac{HTA} system measures tensile and compressive loads using a strain gauge configured in a Wheatstone bridge, connected to an HX711 analog-to-digital converter \cite{Avia:2009}. The acquired data is processed through a software-implemented filter and a PD control algorithm, then output via an MCP4725 digital-to-analog converter \cite{Microchip:2009}.
The analog output is amplified by an LM358 operational amplifier \cite{LM2596:2023} and fed into a transistor-based H-bridge, allowing control of two motors operating simultaneously. Motor velocity is monitored by two LM393 encoders \cite{Rajguru:2017}, closing the feedback loop. A switch and LED provide system power control and status indication.
The Power management is handled by a TPS5430 module \cite{TexasInstruments:2006} providing +/-12V in combination with a LM2596HVS Step-Down Voltage Regulator \cite{TexasInstruments:2023} adding a 5V power source to the System. These two modules are powered by a 12.6V rechargeable LiPo battery and the Hypatia Board has a 5V power bank as recourse.

\begin{figure}[H]
	\centering
	\includegraphics[width=1\linewidth]{../../../Images/Syde/HTA/FlowChartHTA}
	\caption[Flowchart of HTA System]{Flowchart of HTA System}
	\label{fig:FlowChartHTA}
\end{figure}

Having this introduction into the systems and architectures of the demonstrator in mind, the main focus will be on the data transmission throughout the \ac{HTA} system, starting with measuring tensile and compressive loads.

\newpage